\scanpage{550}
\mspage{461}

Je vais résumer ce qui a été obtenu dans un énoncé récapitulatif :

\uline{Théorème} 4.15.9. Soient $C$ une \inserted{petite} catégorie, $\pi$ un cardinal, tel que $\pi\geq\operatorname{card} C$. Les conditions suivantes sont équivalentes (et je conjecture qu'elles sont toujours satisfaites) :

a) $\deleted{\unsure{C|M}}\unsure{\mathcal M|M}\overset{\mathrm{def}}{=}\mathrm{Ind}_\pi C$ est réduit à $\mathcal M_\pi$

a$'$) (Si $C$ est \unsure{Karoubienne|karoubienne}) Tout système inductif $\underline{x}=(x_i)_{i\in I}$ dans $C$, grand devant $\pi$, est essentiellement constant.

a$''$) Pour \unsure{tt} système inductif $(x_i)_{i\in I}$ dans $C$, grand devant $\pi$, il existe un foncteur cofinal $J\mathbin{\unsure{-|\to}}I$, avec $J$ \inserted{ens ordonné} \unsure{grd} devant $\pi$, et tel que le système réindexé \inserted{$(x_j)_{j\in J}$} satisfasse les
conditions suivantes : \unsure{1)|a)} tous les $x_j$ sont égaux à un même élément $x$ de $C$, \unsure{2)|b)} il existe un projecteur $p$ de $x$, tel que l'on ait $\alpha_{j'j}=p$ pour $j<j'$ $(j,j'\in J)$, où $\alpha_{j',j}:x_j\to\unsure{x_{j'}|x_j,|x_j}$ est la flèche de transition.

b) Pour tout $x\in\mathrm{Ob}\,C$, et pour tt monoïde quotient $M'$ de $M=\operatorname{End}_C(x)$, tel que $M'$ soit un monoïde de monomorphismes (i.e. $uv=uv'\Longrightarrow v=v'$ si $u,v,v'\in M'$), et pour tout système
inductif grand devant $\pi$ à valeurs dans $B_{M'}$, (défini par un système d'éléments $(\alpha_{ji})_{\substack{i\leq j\\i,j\in I}}$ d'un ens ordonné $I$ \unsure{grd} devant $\pi$, avec $\alpha_{hj}\alpha_{ji}=\alpha_{hi}$ pour $i\leq j\leq h$), il existe $i_0\in I$ tel
% archive sheet number: 549
